% Pattern: proof/attach/announce. Proofs whose target is named outside a
% plain optional argument: a head that uses a macro of the paper around
% \ref, a head in parentheses inside a font command, a sentence just before
% the proof that announces it, and a subsection head just before it.
% Seen in the gold pilot (set E).
% want lemma-1 statement/lemma 1 status=here labels=l:gap
% want theorem-1 statement/theorem 2 status=here labels=first
% want theorem-2 statement/theorem 3 status=here labels=second
% want theorem-3 statement/theorem 4 status=here labels=third
% want lemma-2 statement/lemma 5 status=none
% want proof-1 proof - proves=lemma-1
% want proof-2 proof - proves=theorem-1
% want proof-3 proof - proves=theorem-2
% want proof-4 proof - proves=theorem-3
\documentclass{article}
\newtheorem{theorem}{Theorem}
\newtheorem{lemma}[theorem]{Lemma}
\newcommand{\refl}[1]{Lemma~\ref{l:#1}}
\begin{document}
\begin{lemma}\label{l:gap}
A helper.
\end{lemma}
\begin{theorem}\label{first}
A statement.
\end{theorem}
\begin{theorem}\label{second}
A second statement.
\end{theorem}
\begin{theorem}\label{third}
A third statement.
\end{theorem}
\begin{lemma}
Another helper.
\end{lemma}
Some text about the lemma.

More text.

\begin{proof}[Proof of \refl{gap}]
A step.
\end{proof}
\begin{proof} \textit{(of Theorem \ref{first})}
Another step.
\end{proof}
We are now ready for the proof of Theorem~\ref{second}.\\
\begin{proof}
A third step.
\end{proof}
\subsection{Proof of Theorem~\ref{third}}
\label{secproof}
\begin{proof}
A fourth step.
\end{proof}
\end{document}
